\section{Functions}

\subsection{Section Overview}
Functions are one of the most fundamental building blocks in C++ and programming in general. They allow us to break down our code into smaller, manageable, and reusable pieces. In this section, we'll cover how to define and call functions, how to pass data to them, how to return values, and more advanced topics like function overloading and recursion.

\subsection{What is a Function?}
A function is a named block of code that performs a specific task. We've already been using functions, such as \texttt{main()}, \texttt{std::cout <<}, and \texttt{strlen()}. Functions help in making programs modular, easier to debug, and promote code reuse.

\subsection{Coding Exercise 22: Using Functions from the cmath Library}
The \texttt{<cmath>} library contains many useful mathematical functions. Write a program that asks the user for a number, then prints its square root (\texttt{sqrt}), cosine (\texttt{cos}), and sine (\texttt{sin}).
\begin{lstlisting}
#include <iostream>
#include <cmath>

int main() {
	double num;
	std::cout << "Enter a number: ";
	std::cin >> num;
	std::cout << "Square root: " << sqrt(num) << std::endl;
	std::cout << "Cosine: " << cos(num) << std::endl;
	std::cout << "Sine: " << sin(num) << std::endl;
	return 0;
}
\end{lstlisting}

\subsection{Function Definition}
A function definition consists of a function header and a function body. The header includes the return type, the function name, and a list of parameters. The body contains the code that the function executes.
\begin{lstlisting}
// return_type function_name (parameter_list) {
//     function_body
// }

int add(int a, int b) {
	return a + b;
}
\end{lstlisting}

\subsection{Function Prototypes}
A function prototype is a declaration of a function that tells the compiler about the function's name, return type, and parameters. It allows you to call a function before you've defined it. Prototypes are often placed in header files.
\begin{lstlisting}
// function prototype
int add(int a, int b);

int main() {
	int sum = add(5, 3); // works because of the prototype
	return 0;
}

// function definition
int add(int a, int b) {
	return a + b;
}
\end{lstlisting}

\subsection{Function Parameters and the return Statement}
Parameters are variables in a function's parameter list that get initialized with the values passed in when the function is called. The \texttt{return} statement specifies the value that the function sends back to the caller. A function with a \texttt{void} return type does not return a value.

\subsection{Coding Exercise 23: Functions and Prototypes --- Converting Temperatures}
Write a function that converts Fahrenheit to Celsius. The formula is C = (F - 32) * 5/9. Create a prototype for this function. In \texttt{main}, call the function with a value and print the result.
\begin{lstlisting}
#include <iostream>

// prototype
double fahrenheit_to_celsius(double f);

int main() {
	double fahrenheit = 98.6;
	double celsius = fahrenheit_to_celsius(fahrenheit);
	std::cout << fahrenheit << "F is " << celsius << "C" << std::endl;
	return 0;
}

// definition
double fahrenheit_to_celsius(double f) {
	return (f - 32.0) * 5.0 / 9.0;
}
\end{lstlisting}

\subsection{Default Argument Values}
You can provide default values for function parameters. If a call to the function omits an argument, the default value is used.
\begin{lstlisting}
void print_greeting(std::string name = "Guest", int age = 0) {
	std::cout << "Hello, " << name << "! You are "
	          << age << " years old." << std::endl;
}

int main() {
	print_greeting("Alice", 25); // Hello, Alice! You are 25 years old.
	print_greeting("Bob");        // Hello, Bob! You are 0 years old.
	print_greeting();             // Hello, Guest! You are 0 years old.
}
\end{lstlisting}

\subsection{Coding Exercise 24: Using Default Argument Values --- Grocery List}
Create a function \texttt{print\_grocery\_item} that takes a \texttt{std::string} for the item name and an \texttt{int} for the quantity with a default value of 1.
\begin{lstlisting}
#include <iostream>
#include <string>

void print_grocery_item(std::string item, int quantity = 1) {
	std::cout << "Item: " << item
	          << ", Quantity: " << quantity << std::endl;
}

int main() {
	print_grocery_item("Apples", 5);
	print_grocery_item("Milk");
	return 0;
}
\end{lstlisting}

\subsection{Overloading Functions}
Function overloading is defining multiple functions with the same name but different parameter lists (different number of parameters or different types). The compiler determines which version of the function to call based on the arguments provided.
\begin{lstlisting}
int add(int a, int b) { return a + b; }
double add(double a, double b) { return a + b; }
int add(int a, int b, int c) { return a + b + c; }
\end{lstlisting}

\subsection{Coding Exercise 25: Overloading Functions --- Calculating Area}
Create two overloaded functions named \texttt{calculate\_area}. One should calculate the area of a square (taking one \texttt{double} for the side) and the other should calculate the area of a rectangle (taking two \texttt{double}s for length and width).
\begin{lstlisting}
#include <iostream>

double calculate_area(double side) {
	return side * side;
}
double calculate_area(double length, double width) {
	return length * width;
}

int main() {
	std::cout << "Area of square: "
	          << calculate_area(10.0) << std::endl;
	std::cout << "Area of rectangle: "
	          << calculate_area(10.0, 20.0) << std::endl;
	return 0;
}
\end{lstlisting}

\subsection{Passing Arrays to Functions}
When you pass an array to a function, you are actually passing a pointer to the first element. The function does not know the size of the array, so you must pass the size as another argument.
\begin{lstlisting}
void print_array(const int arr[], size_t size) {
	for (size_t i = 0; i < size; ++i) {
		std::cout << arr[i] << " ";
	}
	std::cout << std::endl;
}

int main() {
	int my_array[] {1, 2, 3, 4, 5};
	print_array(my_array, 5);
}
\end{lstlisting}
\begin{remark}
Note that any changes made to the array inside the function will affect the original array. The \texttt{const} keyword can be used to prevent modifications.
\end{remark}

\subsection{Coding Exercise 26: Passing Arrays to Functions --- Print a Guest List}
Write a function that takes an array of \texttt{std::string} and its size, and prints each name in the array.
\begin{lstlisting}
#include <iostream>
#include <string>

void print_guest_list(const std::string guests[],
                      size_t size) {
	for (size_t i = 0; i < size; ++i) {
		std::cout << guests[i] << std::endl;
	}
}

int main() {
	std::string guest_list[] {"Alice", "Bob", "Charlie"};
	print_guest_list(guest_list, 3);
	return 0;
}
\end{lstlisting}

\subsection{Pass by Reference}
By default, C++ passes arguments by value (a copy is made). Pass by reference allows a function to access and modify the original argument. This is done by using a reference parameter (\texttt{\&}).
\begin{lstlisting}
void double_it(int &num) {
	num *= 2;
}

int main() {
	int x = 10;
	double_it(x);
	std::cout << "x is now " << x << std::endl; // x is now 20
}
\end{lstlisting}

\subsection{Coding Exercise 27: Using Pass by Reference --- Print a Guest List}
Modify the previous \texttt{print\_guest\_list} function. Create another function \texttt{add\_guest} that takes a \texttt{std::vector<std::string>} by reference and a \texttt{std::string} for the new guest, and adds the guest to the vector.
\begin{lstlisting}
#include <iostream>
#include <vector>
#include <string>

void add_guest(std::vector<std::string> &guests,
               std::string new_guest) {
	guests.push_back(new_guest);
}

void print_guests(
		const std::vector<std::string> &guests) {
	for (const auto &guest : guests) {
		std::cout << guest << std::endl;
	}
}

int main() {
	std::vector<std::string> guest_list {"Alice", "Bob"};
	add_guest(guest_list, "Charlie");
	print_guests(guest_list);
	return 0;
}
\end{lstlisting}

\subsection{Scope Rules}
Scope determines the visibility of an identifier. A local variable, declared inside a function, is only visible within that function. A global variable, declared outside all functions, is visible throughout the program file. It's best to minimize the use of global variables.

\subsection{How do Function Calls Work?}
When a function is called, the program pushes the function's parameters and the return address onto the call stack. The function executes, and when it returns, the return value is sent back and the stack is cleaned up.

\subsection{inline Functions}
For very short functions, you can suggest to the compiler that it perform \texttt{inline} expansion. This copies the function's code directly into the calling location, which can avoid the overhead of a function call. This is only a hint to the compiler.
\begin{lstlisting}
inline int add(int a, int b) { return a + b; }
\end{lstlisting}

\subsection{Recursive Functions}
A recursive function is a function that calls itself, either directly or indirectly. A recursive function must have a base case (a condition that stops the recursion) to avoid an infinite loop.
\begin{lstlisting}
// factorial function
unsigned long long factorial(unsigned long long n) {
	if (n == 0)
		return 1; // base case
	return n * factorial(n - 1); // recursive step
}
\end{lstlisting}

\subsection{Coding Exercise 28: Implementing a Recursive Function --- Sum of Digits}
Write a recursive function to find the sum of the digits of a positive integer.
\begin{lstlisting}
#include <iostream>

int sum_of_digits(int n) {
	if (n == 0) {
		return 0;
	}
	return (n % 10) + sum_of_digits(n / 10);
}

int main() {
	std::cout << "Sum of digits of 123 is "
	          << sum_of_digits(123) << std::endl;
	return 0;
}
\end{lstlisting}

\subsection{Coding Exercise 29: Implementing a Recursive Function --- Save a Penny}
Calculate how much money you would have if you saved a penny on day 1, and doubled the amount you save each day for 30 days. Implement this with a recursive function.
\begin{lstlisting}
#include <iostream>
#include <cmath>

double calculate_savings(int days) {
	if (days == 1) {
		return 0.01;
	}
	return pow(2, days - 1) * 0.01
	       + calculate_savings(days - 1);
}

int main() {
	std::cout.precision(10);
	std::cout << "Total savings after 30 days: $"
	          << calculate_savings(30) << std::endl;
	return 0;
}
\end{lstlisting}

\subsection{Section Challenge}
Write a program that uses functions to provide a menu-driven interface similar to the one in the ``Controlling Program Flow'' section. Create separate functions for printing the menu, adding a number, calculating the mean, etc.

\subsection{Section Challenge --- Solution}
\begin{lstlisting}
#include <iostream>
#include <vector>
#include <string>

void print_menu();
char get_choice();
void print_numbers(const std::vector<int> &numbers);
void add_number(std::vector<int> &numbers);
void display_mean(const std::vector<int> &numbers);
void display_smallest(const std::vector<int> &numbers);
void display_largest(const std::vector<int> &numbers);

int main() {
	std::vector<int> numbers;
	char choice;
	do {
		print_menu();
		choice = get_choice();
		switch (choice) {
			case 'P': print_numbers(numbers); break;
			case 'A': add_number(numbers); break;
			case 'M': display_mean(numbers); break;
			case 'S': display_smallest(numbers); break;
			case 'L': display_largest(numbers); break;
			case 'Q':
				std::cout << "Goodbye!" << std::endl;
				break;
			default:
				std::cout << "Unknown selection, "
				          << "please try again"
				          << std::endl;
		}
	} while (choice != 'Q');
	return 0;
}

void print_menu() {
	std::cout << "\nP - Print numbers\n"
	          << "A - Add a number\n"
	          << "M - Display mean\n"
	          << "S - Display smallest\n"
	          << "L - Display largest\n"
	          << "Q - Quit\n"
	          << "\nEnter your choice: ";
}

char get_choice() {
	char c;
	std::cin >> c;
	return toupper(c);
}

void print_numbers(
		const std::vector<int> &numbers) {
	if (numbers.empty()) {
		std::cout << "[] - the list is empty"
		          << std::endl;
	} else {
		std::cout << "[ ";
		for (int num : numbers)
			std::cout << num << " ";
		std::cout << "]" << std::endl;
	}
}

void add_number(std::vector<int> &numbers) {
	int num;
	std::cout << "Enter an integer to add: ";
	std::cin >> num;
	numbers.push_back(num);
	std::cout << num << " added." << std::endl;
}

void display_mean(
		const std::vector<int> &numbers) {
	if (numbers.empty()) {
		std::cout << "Unable to calculate mean"
		          << " - no data" << std::endl;
	} else {
		int sum = 0;
		for (int num : numbers) sum += num;
		std::cout << "Mean: "
		          << static_cast<double>(sum)
		             / numbers.size()
		          << std::endl;
	}
}

void display_smallest(
		const std::vector<int> &numbers) {
	if (numbers.empty()) {
		std::cout << "Unable to determine smallest"
		          << " - no data" << std::endl;
	} else {
		int smallest = numbers[0];
		for (int num : numbers)
			if (num < smallest) smallest = num;
		std::cout << "Smallest number: "
		          << smallest << std::endl;
	}
}

void display_largest(
		const std::vector<int> &numbers) {
	if (numbers.empty()) {
		std::cout << "Unable to determine largest"
		          << " - no data" << std::endl;
	} else {
		int largest = numbers[0];
		for (int num : numbers)
			if (num > largest) largest = num;
		std::cout << "Largest number: "
		          << largest << std::endl;
	}
}
\end{lstlisting}

\subsection{Quiz 8: Section 11 Quiz}
\begin{enumerate}
    \item What is a function prototype?
    \begin{enumerate}
        \item[a)] The body of a function.
        \item[b)] A declaration of a function that specifies its name, return type, and parameters.
        \item[c)] A way to call a function.
        \item[d)] Another name for the \texttt{main} function.
    \end{enumerate}

    \item What is function overloading?
    \begin{enumerate}
        \item[a)] A function that calls itself.
        \item[b)] Creating multiple functions with the same name but different parameter lists.
        \item[c)] Providing default values for function parameters.
        \item[d)] Passing a function as an argument to another function.
    \end{enumerate}

    \item When you pass an argument by reference, what is actually passed to the function?
    \begin{enumerate}
        \item[a)] A copy of the argument's value.
        \item[b)] The memory address of the argument.
        \item[c)] The argument's type.
        \item[d)] Nothing is passed.
    \end{enumerate}
\end{enumerate}

\textbf{Answers:} 1. (b), 2. (b), 3. (b)
